\scanpage{7}

- d -

\uline{Commentaires}

Le rédacteur a suivi aussi fidèlement que possible les instructions de La Tribu n$^{\mathrm{o}}$ 57. Le plan n'a été modifié qu'en ce qui concerne les remords-fourre-tout du début, que le rédacteur a jugé préférable de grouper en 3 paragraphes préliminaires au lieu d'un, le 3e étant consacré aux lemmes sur les graphes qui traînaient un peu partout, et dont en fait on a besoin dès le début, (p. ex. pour le critère pour qu'une matrice de Coxeter soit connexe). Conformément aux injonctions de La Tribu, le rédacteur a suivi la terminologie de Berge pour les graphes ; mais si les termes de ``sommet'' et ``arête'' sont classiques, le rédacteur n'aime guère ``sommet pendant'' qui a l'air grotesque, et préférerait ``point d'arrêt'', qui irait assez bien avec ``point de ramification''.

Autre question de terminologie : le Congrès avait adopté d'enthousiasme une belle terminologie symétrique : ``poids'', ``poids radiciels'', ``copoids'', ``copoids radiciels''. Le malheur est qu'avec la définition adoptée, un ``copoids radiciel'' \uline{n'est pas} nécessairement un ``copoids'', ce qui est grammaticalement intolérable. Le rédacteur voit trois solutions possibles :

1) Prendre une autre terminologie pour ``poids radiciels'' et ``copoids radiciels'' ; c'est la solution adoptée dans la rédaction, qui dit respectivement ``vocables'' et ``covocables'', (dactylo, attention, ce n'est \uline{pas} ``convocables''!). Le rédacteur ne tient pas spécialement à ces termes ; on pourrait aussi dire ``noeuds'' et ``conoeuds'' bien que ce dernier mot ait l'air un peu ... ce que vous pensez.

2) On ne peut naturellement changer le sens reçu de ``racines'' et ``poids'' ; mais ``copoids'' est beaucoup plus récent, (Chevalley, je crois), et pas encore très répandu ; on pourrait donc intervertir les sens de ``copoids'' et ``copoids radiciels'', les poids étant donc associés aux copoids radiciels, (et non plus aux copoids), et les copoids radiciels étant engendrés par les coracines, (ce qui est assez satisfaisant).

3) N'utiliser aucune terminologie pour ``poids radiciels'' et ``copoids radiciels'' ; l'expérience, (tout au moins dans la présente rédaction), montre qu'on s'en sert en fait très peu et qu'on s'en passe fort bien.

\mspage{6}
